\chapter{Photorefractive Keratectomy (PRK)}

\section*{Overview}
PRK is a laser vision-correction procedure that reshapes the cornea without creating a corneal flap. It is an elective option for selected people with refractive errors, with slower surface healing than LASIK and a different recovery experience.

\section*{Why It Is Done}
It is an elective alternative to LASIK for suitable patients, including some with thin corneas or occupations where flap injury is a concern.

\section*{How It Is Performed}
The surface corneal epithelium is removed, an excimer laser reshapes the cornea, and a bandage contact lens is placed while the surface heals.

\section*{Risks and Recovery}
Pain during early healing, haze, infection, dry eye, glare, residual refractive error, and regression are possible. Visual recovery is slower than after LASIK and requires prescribed drops and follow-up.

\section*{Outcome and Follow-up}
The expected outcome depends on the reason for the procedure, the person's health, and whether additional treatment is needed. The clinician reviews the result with the patient, explains any next steps, and sets follow-up according to the findings and recovery.

\section*{Contraindications and Precautions}
Active eye infection, unstable refraction, significant corneal disease, severe uncontrolled dry eye, or inadequate corneal tissue may make PRK unsuitable. Pregnancy and breastfeeding can temporarily change refraction; the eye surgeon reviews healing risk, medicines, and expectations.

\section*{When to Seek Medical Care}
Follow the procedure team's specific instructions. Seek urgent medical assessment for severe or worsening pain, uncontrolled bleeding, fever or signs of infection, trouble breathing, chest pain, fainting, new weakness or confusion, inability to urinate, or any rapidly worsening symptom. The appropriate warning signs depend on the procedure and the patient's medical condition.

\section*{References}
\begin{enumerate}
  \item National Eye Institute. Refractive Surgery. 2025.
  \item Current Medical Diagnosis and Treatment. McGraw-Hill; 2025.
\end{enumerate}
\clearpage
